\documentclass[11pt]{article}
\usepackage[margin=1in]{geometry}
\usepackage{listings,xcolor,hyperref,booktabs,longtable,amsmath}
\lstset{basicstyle=\ttfamily\small,breaklines=true,frame=single,numbers=left,numberstyle=\tiny\color{gray}}
\title{Sentence Embedding, Chunking \& RAG Pipeline:\\Quality-First Retrieval-Augmented Generation}
\author{SAP OSS Technical Documentation}
\date{March 2026}

\begin{document}
\maketitle
\tableofcontents
\newpage

% ============================================================================
\section{Introduction}
% ============================================================================

This report documents the sentence embedding, text chunking, and Retrieval-Augmented
Generation (RAG) pipeline implemented across the SAP OSS codebase. The system spans
three primary codebases:

\begin{enumerate}
  \item \textbf{LangChain Integration for SAP HANA Cloud} --- MCP server with
        sentence-aware text splitting, KùzuDB graph-RAG, and Python LangChain
        vector store integration.
  \item \textbf{CAP LLM Plugin} --- TypeScript CDS service providing embedding
        generation, similarity search, RAG pipeline orchestration, and response
        quality assessment.
  \item \textbf{OData Vocabularies} --- Graph-RAG store for vocabulary term
        relationships and annotation targets.
\end{enumerate}

The architecture follows a quality-first design philosophy: every stage---from
chunking through retrieval through generation---includes explicit quality controls,
observability spans, and graceful degradation paths.

% ============================================================================
\section{Sentence-Aware Recursive Splitting}
\label{sec:recursive-split}
% ============================================================================

\noindent\textbf{Source:}
\texttt{src/data/langchain-integration-for-sap-hana-cloud-main/mcp\_server/server.py},
lines 72--128.

The \texttt{\_recursive\_split()} function implements a hierarchical text splitter
that attempts coarse-grained separators first and recursively falls back to
finer-grained ones. The separator hierarchy is:

\begin{center}
\begin{tabular}{cl}
\toprule
\textbf{Priority} & \textbf{Separator} \\
\midrule
1 & \verb|"\n\n"| (paragraph boundary) \\
2 & \verb|"\n"| (line boundary) \\
3 & \verb|". "| (sentence-ending period) \\
4 & \verb|"! "| (exclamation) \\
5 & \verb|"? "| (question mark) \\
6 & \verb|"; "| (semicolon) \\
7 & \verb|", "| (comma) \\
8 & \verb|" "| (word boundary) \\
9 & \verb|""| (character-level, base case) \\
\bottomrule
\end{tabular}
\end{center}

\subsection{Function Signature and Early Return}

\begin{lstlisting}[language=Python,caption={\_recursive\_split() --- lines 72--79}]
def _recursive_split(text: str, separators: list,
                     chunk_size: int, overlap: int) -> list:
    """Recursively split text using decreasing separator granularity.

    Separator is preserved at the end of each part (e.g. splitting
    on ". " keeps the period+space attached to the preceding chunk)
    so no text is lost.
    """
    if len(text) <= chunk_size:
        return [text] if text.strip() else []
\end{lstlisting}

The early return at line 78 ensures that text already fitting within
\texttt{chunk\_size} is returned as a single chunk (or discarded if whitespace-only).

\subsection{Separator Selection}

\begin{lstlisting}[language=Python,caption={Separator selection --- lines 81--82}]
    sep = separators[0] if separators else ""
    remaining_seps = separators[1:] if len(separators) > 1 else [""]
\end{lstlisting}

The current separator is popped from the front of the list. The remaining
separators are passed to recursive calls for finer-grained splitting.

\subsection{Base Case: Character-Level Splitting}

When the separator list is exhausted (empty string), the function falls back to
hard character splitting with overlap:

\begin{lstlisting}[language=Python,caption={Character-level base case --- lines 84--94}]
    if not sep:
        # Base case: hard character split with overlap
        chunks = []
        step = max(chunk_size - overlap, 1)
        for i in range(0, len(text), step):
            chunk = text[i:i + chunk_size]
            if chunk.strip():
                chunks.append(chunk)
            if i + chunk_size >= len(text):
                break
        return chunks
\end{lstlisting}

The step size is \texttt{chunk\_size - overlap}, ensuring that consecutive chunks
share \texttt{overlap} characters. The guard \texttt{max(..., 1)} prevents infinite
loops when overlap equals or exceeds the chunk size.

\subsection{Separator Preservation}

A critical design decision is that separators are \emph{re-attached} to the
preceding part, ensuring no text is lost during splitting:

\begin{lstlisting}[language=Python,caption={Separator preservation --- lines 96--101}]
    # Split then re-attach separator to each part so no text is lost.
    # "A. B. C" split on ". " -> raw ["A", "B", "C"] -> ["A. ", "B. ", "C"]
    raw_parts = text.split(sep)
    parts = []
    for i, p in enumerate(raw_parts):
        parts.append(p + sep if i < len(raw_parts) - 1 else p)
\end{lstlisting}

\subsection{Chunk Assembly with Overlap}

The greedy accumulation loop (lines 103--127) assembles chunks:

\begin{lstlisting}[language=Python,caption={Chunk assembly --- lines 103--128}]
    chunks = []
    current = ""

    for part in parts:
        candidate = current + part
        if len(candidate) <= chunk_size:
            current = candidate
        else:
            if current.strip():
                chunks.append(current)
            # If this single part exceeds chunk_size, recurse with
            # finer separator
            if len(part) > chunk_size:
                chunks.extend(_recursive_split(
                    part, remaining_seps, chunk_size, overlap))
                current = ""
            else:
                # Start new chunk with overlap from previous
                if overlap > 0 and current:
                    tail = current[-overlap:]
                    current = tail + part if tail.strip() else part
                else:
                    current = part

    if current.strip():
        chunks.append(current)

    return chunks
\end{lstlisting}

Key behaviors:
\begin{itemize}
  \item Parts that fit within \texttt{chunk\_size} are greedily accumulated.
  \item When a single part exceeds \texttt{chunk\_size}, it is recursively split
        using the next finer separator in the hierarchy.
  \item Overlap is implemented by copying the last \texttt{overlap} characters
        from the previous chunk as a prefix to the new chunk.
  \item Whitespace-only chunks are discarded at every stage.
\end{itemize}

% ============================================================================
\section{MCP Text Splitter Tool}
\label{sec:mcp-splitter}
% ============================================================================

\noindent\textbf{Source:}
\texttt{src/data/langchain-integration-for-sap-hana-cloud-main/mcp\_server/server.py},
lines 412--424 (tool schema) and lines 785--802 (handler).

\subsection{Tool Schema}

The text splitter is exposed as an MCP tool named \texttt{langchain\_split\_text}:

\begin{lstlisting}[language=Python,caption={Tool registration --- lines 412--424}]
self.tools["langchain_split_text"] = {
    "name": "langchain_split_text",
    "description": "Split text using sentence-aware recursive splitter "
                   "(paragraph -> sentence -> word boundaries)",
    "inputSchema": {
        "type": "object",
        "properties": {
            "text": {"type": "string",
                     "description": "Text to split"},
            "chunk_size": {"type": "number",
                "description": "Target chunk size in characters "
                               "(default 1000, max 4000)"},
            "chunk_overlap": {"type": "number",
                "description": "Overlap between chunks in "
                               "characters (default 200)"},
        },
        "required": ["text"],
    },
}
\end{lstlisting}

\subsection{Handler Implementation}

\begin{lstlisting}[language=Python,caption={Split text handler --- lines 785--802}]
def _handle_langchain_split_text(self, args: dict) -> dict:
    text = args.get("text", "")
    if not text:
        return {"chunks": 0, "chunk_size": 0,
                "overlap": 0, "texts": []}

    chunk_size = clamp_int(
        args.get("chunk_size", 1000), 1000, 1, MAX_CHUNK_SIZE)
    overlap = clamp_int(
        args.get("chunk_overlap", 200), 200, 0, chunk_size - 1)

    # Sentence-aware recursive splitting:
    # paragraph -> sentence -> word -> char
    separators = ["\n\n", "\n", ". ", "! ", "? ",
                  "; ", ", ", " ", ""]
    chunks = _recursive_split(text, separators,
                              chunk_size, overlap)

    return {
        "chunks": len(chunks),
        "chunk_size": chunk_size,
        "overlap": overlap,
        "texts": chunks,
    }
\end{lstlisting}

The handler enforces safe defaults via \texttt{clamp\_int()}: chunk size is clamped
to $[1, 4000]$ (governed by \texttt{MAX\_CHUNK\_SIZE} at line 46), and overlap is
clamped to $[0, \texttt{chunk\_size} - 1]$ to prevent degenerate cases. The return
payload includes both the chunk metadata (\texttt{chunks}, \texttt{chunk\_size},
\texttt{overlap}) and the actual split texts.

% ============================================================================
\section{Embedding Model Registry}
\label{sec:model-registry}
% ============================================================================

\noindent\textbf{Source:}
\texttt{src/data/cap-llm-plugin-main/lib/validation-utils.js}, lines 69--112.

\subsection{Dimension Map}

The \texttt{EMBEDDING\_MODEL\_DIMENSIONS} constant maps 15 known embedding models
to their output dimensionality:

\begin{lstlisting}[language=Java,caption={Model registry --- lines 73--94}]
const EMBEDDING_MODEL_DIMENSIONS = {
  // OpenAI
  "text-embedding-ada-002": 1536,
  "text-embedding-3-small": 1536,
  "text-embedding-3-large": 3072,
  // OpenAI with custom dimensions
  "text-embedding-3-small-256": 256,
  "text-embedding-3-small-512": 512,
  // Cohere
  "embed-english-v3.0": 1024,
  "embed-multilingual-v3.0": 1024,
  "embed-english-light-v3.0": 384,
  "embed-multilingual-light-v3.0": 384,
  // Voyage
  "voyage-large-2": 1536,
  "voyage-code-2": 1536,
  "voyage-2": 1024,
  // BGE (common open-source, used via SAP AI Core)
  "bge-large-en-v1.5": 1024,
  "bge-base-en-v1.5": 768,
  "bge-small-en-v1.5": 384,
};
\end{lstlisting}

\begin{longtable}{llr}
\toprule
\textbf{Provider} & \textbf{Model} & \textbf{Dimensions} \\
\midrule
\endhead
OpenAI & text-embedding-ada-002 & 1536 \\
OpenAI & text-embedding-3-small & 1536 \\
OpenAI & text-embedding-3-large & 3072 \\
OpenAI & text-embedding-3-small-256 & 256 \\
OpenAI & text-embedding-3-small-512 & 512 \\
Cohere & embed-english-v3.0 & 1024 \\
Cohere & embed-multilingual-v3.0 & 1024 \\
Cohere & embed-english-light-v3.0 & 384 \\
Cohere & embed-multilingual-light-v3.0 & 384 \\
Voyage & voyage-large-2 & 1536 \\
Voyage & voyage-code-2 & 1536 \\
Voyage & voyage-2 & 1024 \\
BGE & bge-large-en-v1.5 & 1024 \\
BGE & bge-base-en-v1.5 & 768 \\
BGE & bge-small-en-v1.5 & 384 \\
\bottomrule
\end{longtable}

\subsection{Dimension Validation}

\begin{lstlisting}[language=Java,caption={validateEmbeddingDimensions() --- lines 104--112}]
function validateEmbeddingDimensions(embedding, modelName) {
  const expectedDims = EMBEDDING_MODEL_DIMENSIONS[modelName];
  if (expectedDims && embedding.length !== expectedDims) {
    throw new Error(
      `Embedding dimension mismatch: model "${modelName}" ` +
      `produces ${expectedDims}-dim vectors, ` +
      `but received ${embedding.length} dimensions. ` +
      `Ensure the same model is used for indexing and querying.`
    );
  }
}
\end{lstlisting}

The function exhibits two distinct behaviors:
\begin{enumerate}
  \item \textbf{Known model, wrong dimensions:} Throws an \texttt{Error} with a
        diagnostic message identifying the expected vs.\ actual dimension count.
  \item \textbf{Unknown model:} Passes silently (no-op), allowing custom or
        newly released models to work without code changes.
\end{enumerate}

This is invoked in the RAG pipeline at
\texttt{cap-llm-plugin.ts} line 466:
\begin{lstlisting}[language=Java]
validateEmbeddingDimensions(queryEmbedding, embeddingConfig.modelName);
\end{lstlisting}

% ============================================================================
\section{Vector Embedding Generation}
\label{sec:embedding-gen}
% ============================================================================

\noindent\textbf{Source:}
\texttt{src/data/cap-llm-plugin-main/srv/cap-llm-plugin.ts}, lines 244--296.

\subsection{getEmbeddingWithConfig()}

The \texttt{getEmbeddingWithConfig()} method generates vector embeddings using the
SAP AI SDK's \texttt{OrchestrationEmbeddingClient}:

\begin{lstlisting}[language=Java,caption={getEmbeddingWithConfig() --- lines 244--296}]
async getEmbeddingWithConfig(
  config: EmbeddingConfig, input: unknown
): Promise<unknown> {
  const span = getTracer().startSpan(
    "cap-llm-plugin.getEmbeddingWithConfig");
  span.setAttribute("llm.embedding.model",
    config?.modelName ?? "");
  span.setAttribute("llm.resource_group",
    config?.resourceGroup ?? "");

  try {
    if (!config?.modelName) {
      throw new EmbeddingError(
        `The config is missing the parameter: "modelName".`,
        "EMBEDDING_CONFIG_INVALID",
        { missingField: "modelName" });
    }
    if (!config?.resourceGroup) {
      throw new EmbeddingError(
        `The config is missing the parameter: "resourceGroup".`,
        "EMBEDDING_CONFIG_INVALID",
        { missingField: "resourceGroup" });
    }

    const { OrchestrationEmbeddingClient } =
      await import("@sap-ai-sdk/orchestration");

    const client = new OrchestrationEmbeddingClient(
      { embeddings: { model: { name: config.modelName as any }}},
      { resourceGroup: config.resourceGroup }
    );

    const response = await client.embed(
      { input: input as string | string[] },
      {
        middleware: [
          createOtelMiddleware({
            endpoint: "/embeddings",
            resourceGroup: config.resourceGroup,
          }),
        ],
      }
    );
    span.addEvent("embedding_response_received");
    span.setStatus({ code: SpanStatusCode.OK });
    return response;
  } catch (e) { /* ... error handling ... */ }
  finally { span.end(); }
}
\end{lstlisting}

Key architectural decisions:
\begin{itemize}
  \item The \texttt{@sap-ai-sdk/orchestration} import is \emph{dynamic}
        (\texttt{await import(...)}) to avoid hard compile-time dependencies.
  \item OTEL (OpenTelemetry) middleware is injected into the SDK call via
        \texttt{createOtelMiddleware()}, enabling distributed tracing across
        the embedding request.
  \item Errors are wrapped in typed \texttt{EmbeddingError} with structured
        error codes (\texttt{EMBEDDING\_CONFIG\_INVALID},
        \texttt{EMBEDDING\_REQUEST\_FAILED}).
\end{itemize}

% ============================================================================
\section{Similarity Search}
\label{sec:similarity-search}
% ============================================================================

\noindent\textbf{Source:}
\texttt{src/data/cap-llm-plugin-main/srv/cap-llm-plugin.ts}, lines 657--762.

The similarity search implementation follows a dual-path strategy: a preferred
SDK-based path and a raw SQL fallback.

\subsection{Input Validation}

All inputs are validated before SQL construction (lines 672--686):

\begin{lstlisting}[language=Java,caption={Input validation --- lines 672--686}]
validateSqlIdentifier(tableName, "tableName");
validateSqlIdentifier(embeddingColumnName, "embeddingColumnName");
validateSqlIdentifier(contentColumn, "contentColumn");
validatePositiveInteger(topK, "topK", 10000);
validateEmbeddingVector(embedding);

const validAlgorithms = ["COSINE_SIMILARITY", "L2DISTANCE"];
if (!validAlgorithms.includes(algoName)) {
  throw new InvalidSimilaritySearchAlgoNameError(
    `Invalid algorithm name: ${algoName}. ` +
    `Currently only COSINE_SIMILARITY and L2DISTANCE ` +
    `are accepted.`, 400);
}
\end{lstlisting}

The \texttt{validateSqlIdentifier()} function (validation-utils.js, lines 19--29)
enforces the regex \texttt{/\^{}[A-Za-z\_][A-Za-z0-9\_-]*\$/}, preventing SQL
injection in identifier positions. The \texttt{validateEmbeddingVector()} function
(lines 54--66) ensures the embedding is a non-empty array of finite numbers.

\subsection{Preferred Path: @sap-ai-sdk/hana-vector}

\begin{lstlisting}[language=Java,caption={HANAVectorStore SDK path --- lines 690--726}]
try {
  const { createHANAClientFromEnv, createHANAVectorStore }
    = await import("@sap-ai-sdk/hana-vector" as string);
  const hanaClient = createHANAClientFromEnv();
  await hanaClient.init();
  const metricMap: Record<string, "COSINE" | "EUCLIDEAN"> = {
    COSINE_SIMILARITY: "COSINE",
    L2DISTANCE: "EUCLIDEAN",
  };
  const vectorStore = createHANAVectorStore(hanaClient, {
    tableName,
    embeddingColumn: embeddingColumnName,
    contentColumn,
    embeddingDimensions: embedding.length,
  });
  const hits = await vectorStore.similaritySearch(embedding, {
    k: topK,
    metric: metricMap[algoName] ?? "COSINE",
  });
  sdkResult = hits.map((h: any) => ({
    PAGE_CONTENT: h.content ?? h.metadata?.content ?? "",
    SCORE: h.score ?? 0,
    ...h,
  })) as SimilaritySearchResult[];
  await hanaClient.disconnect?.();
} catch (sdkErr) {
  LOG.debug("@sap-ai-sdk/hana-vector not available, "
    + "falling back to raw SQL:",
    (sdkErr as Error).message);
  sdkResult = undefined;
}
\end{lstlisting}

\subsection{Fallback Path: Raw SQL}

When the SDK is unavailable, the system falls back to constructing a raw SQL
query with double-quoted identifiers:

\begin{lstlisting}[language=Java,caption={Raw SQL fallback --- lines 728--748}]
let sortDirection = "DESC";
if ("L2DISTANCE" === algoName) {
  sortDirection = "ASC";
}

const embeddingStr = "'[" + embedding.join(",") + "]'";

const selectStmt =
  `SELECT TOP ${Number(topK)} *,` +
  `TO_NVARCHAR("${contentColumn}") as PAGE_CONTENT,` +
  `${algoName}("${embeddingColumnName}", ` +
  `TO_REAL_VECTOR(${embeddingStr})) as SCORE ` +
  `FROM "${tableName}" ORDER BY SCORE ${sortDirection}`;
const db = (await cds.connect.to("db")) as any;
const result = await db.run(selectStmt);
\end{lstlisting}

SQL injection is prevented by:
\begin{enumerate}
  \item \texttt{validateSqlIdentifier()} on all identifier positions.
  \item \texttt{validateEmbeddingVector()} ensuring only finite numbers in the
        embedding string.
  \item \texttt{Number(topK)} coercion for the TOP clause.
  \item Double-quoted identifiers for table and column names.
\end{enumerate}

% ============================================================================
\section{RAG Pipeline Quality Controls}
\label{sec:rag-quality}
% ============================================================================

\noindent\textbf{Source:}
\texttt{src/data/cap-llm-plugin-main/srv/cap-llm-plugin.ts},
\texttt{getRagResponseWithConfig()} at lines 436--597.

The RAG pipeline implements six quality control mechanisms.

\subsection{Similarity Score Threshold Filtering}

After retrieving similarity results, low-quality matches are filtered based on
the algorithm type (lines 479--487):

\begin{lstlisting}[language=Java,caption={Score threshold filtering --- lines 479--487}]
// Fix 2: Filter by minimum similarity score
// Cosine: higher is better (default threshold 0.3).
// L2: lower is better (default threshold 1.5).
const allResults =
  similaritySearchResults as SimilaritySearchResult[];
const isCosineSimilarity = algoName === "COSINE_SIMILARITY";
const minScore = (chatParams?.minSimilarityScore as number)
  ?? (isCosineSimilarity ? 0.3 : undefined);
const maxL2Distance =
  (chatParams?.maxL2Distance as number) ?? 1.5;
const filteredResults = isCosineSimilarity
  ? allResults.filter(r => r.SCORE >= (minScore as number))
  : allResults.filter(r => r.SCORE <= maxL2Distance);
\end{lstlisting}

Default thresholds:
\begin{itemize}
  \item \textbf{COSINE\_SIMILARITY:} $\text{score} \geq 0.3$
  \item \textbf{L2DISTANCE:} $\text{score} \leq 1.5$
\end{itemize}

Both thresholds are overridable via \texttt{chatParams.minSimilarityScore} and
\texttt{chatParams.maxL2Distance}.

\subsection{Context Window Budgeting}

The context budget is calculated with overhead awareness (lines 502--530):

\begin{lstlisting}[language=Java,caption={Context budget calculation --- lines 502--530}]
const maxTotalContextTokens =
  (chatParams?.maxContextTokens as number) ?? 3000;
const instructionChars =
  chatInstruction.length + "\n\nRetrieved context:\n".length;
const queryChars = input.length;
const historyChars = context
  ? context.reduce((sum: number, m: any) =>
      sum + ((m?.content as string)?.length ?? 0), 0)
  : 0;
// Per-document header overhead:
// "[Document N] (relevance: 0.XXX)\n" ~ 40 chars
const perDocOverhead = 40;
const reservedChars = instructionChars + queryChars
  + historyChars + 200; // 200 for message framing
const maxContextChars =
  Math.max(0, maxTotalContextTokens * 4 - reservedChars);
let contextBudget = maxContextChars;
const selectedContent: string[] = [];

for (const result of filteredResults) {
  const content = result.PAGE_CONTENT ?? "";
  if (contextBudget <= perDocOverhead) break;
  const available = contextBudget - perDocOverhead;
  if (content.length <= available) {
    selectedContent.push(content);
    contextBudget -= content.length + perDocOverhead;
  } else {
    // Truncate last chunk to fit budget
    // (at sentence boundary if possible)
    const truncated = content.slice(0, available);
    const lastSentence = truncated.lastIndexOf(". ");
    selectedContent.push(
      lastSentence > 0
        ? truncated.slice(0, lastSentence + 1)
        : truncated);
    break;
  }
}
\end{lstlisting}

The budget formula is:
\[
\text{maxContextChars} = \max\!\big(0,\; \text{maxTokens} \times 4 - (\text{instruction} + \text{query} + \text{history} + 200)\big)
\]
where the factor of 4 is the approximate characters-per-token ratio.

\subsection{Zero-Result Fallback Prompt}

When no documents pass the similarity threshold (lines 537--544):

\begin{lstlisting}[language=Java,caption={Zero-result fallback --- lines 537--544}]
if (selectedContent.length === 0) {
  span.addEvent("rag_no_relevant_context", {
    "search.total_results": allResults.length,
    "search.filtered_results": filteredResults.length,
  });
  systemPrompt = `${chatInstruction}\n\n` +
    `No sufficiently relevant documents were found ` +
    `in the knowledge base. Answer based on your ` +
    `general knowledge and clearly state when you ` +
    `are not certain.`;
}
\end{lstlisting}

\subsection{Structured Context Injection}

When documents are available, they are injected with numbered headers and
relevance scores (lines 547--551):

\begin{lstlisting}[language=Java,caption={Structured context injection --- lines 547--551}]
const contextBlock = selectedContent.map((content, i) => {
  const score = filteredResults[i]?.SCORE ?? 0;
  return `[Document ${i + 1}] (relevance: ` +
    `${score.toFixed(3)})\n${content}`;
}).join("\n\n");
systemPrompt = `${chatInstruction}\n\n` +
  `Retrieved context:\n${contextBlock}`;
\end{lstlisting}

Each document is tagged as \texttt{[Document N] (relevance: X.XXX)}, providing
the LLM with explicit provenance information.

\subsection{Response Quality Assessment}

After receiving the chat completion, the response is evaluated (lines 571--577):

\begin{lstlisting}[language=Java,caption={Quality assessment --- lines 571--577}]
const quality = assessResponseQuality(
  chatCompletionResp, selectedContent);
span.setAttribute("rag.response.estimated_tokens",
  quality.estimatedTokens);
span.setAttribute("rag.response.grounded",
  quality.usedContext);
if (quality.warnings.length > 0) {
  span.addEvent("quality_warnings",
    { warnings: quality.warnings.join("; ") });
}
\end{lstlisting}

The quality object is attached to the final \texttt{RagResponse} (line 583).

\subsection{Response Payload Structure}

The RAG pipeline returns a structured \texttt{RagResponse} (lines 580--584):

\begin{lstlisting}[language=Java,caption={RagResponse interface --- lines 85--95}]
export interface RagResponse {
  completion: unknown;
  additionalContents: SimilaritySearchResult[];
  quality?: {
    hasContent: boolean;
    contentLength: number;
    estimatedTokens: number;
    usedContext: boolean;
    warnings: string[];
  };
}
\end{lstlisting}

% ============================================================================
\section{Response Quality Assessment}
\label{sec:quality-assessment}
% ============================================================================

\noindent\textbf{Source:}
\texttt{src/data/cap-llm-plugin-main/lib/validation-utils.js}, lines 122--183.

The \texttt{assessResponseQuality()} function performs a grounding check to
determine whether the LLM response is anchored in the provided context.

\subsection{Content Extraction}

\begin{lstlisting}[language=Java,caption={Content extraction --- lines 131--133}]
const content = typeof response === "string"
  ? response
  : response?.getContent?.()
    ?? response?.choices?.[0]?.message?.content ?? "";
\end{lstlisting}

The function supports three response formats: raw strings, SDK responses with
\texttt{getContent()}, and OpenAI-format responses with
\texttt{choices[0].message.content}.

\subsection{Grounding Check: Bigram + Keyword Overlap}

The grounding check uses two complementary signals (lines 144--179):

\begin{lstlisting}[language=Java,caption={Grounding check --- lines 144--179}]
const normalize = (s) =>
  s.toLowerCase().replace(/[^\w\s]/g, "");
const contextNorm = normalize(context.join(" "));
const responseNorm = normalize(content);

// Keyword overlap (words > 4 chars, filters stopwords)
const contextWords = new Set(
  contextNorm.split(/\s+/).filter(w => w.length > 4));
const responseWords =
  responseNorm.split(/\s+/).filter(w => w.length > 4);
let wordOverlap = 0;
for (const word of responseWords) {
  if (contextWords.has(word)) wordOverlap++;
}

// Bigram overlap (multi-word phrases, more robust)
const toBigrams = (words) => {
  const bigrams = new Set();
  for (let i = 0; i < words.length - 1; i++) {
    bigrams.add(words[i] + " " + words[i + 1]);
  }
  return bigrams;
};
const contextBigrams = toBigrams(
  contextNorm.split(/\s+/).filter(w => w.length > 2));
const responseBigramList =
  responseNorm.split(/\s+/).filter(w => w.length > 2);
let bigramOverlap = 0;
for (let i = 0; i < responseBigramList.length - 1; i++) {
  if (contextBigrams.has(responseBigramList[i]
      + " " + responseBigramList[i + 1])) {
    bigramOverlap++;
  }
}

// Grounded if: significant keyword overlap
// OR any bigram matches (stronger signal)
assessment.usedContext =
  wordOverlap > 3 || bigramOverlap > 1;
assessment.groundingDetail =
  { wordOverlap, bigramOverlap };
\end{lstlisting}

The grounding decision uses two thresholds:
\begin{itemize}
  \item \textbf{Keyword overlap} $> 3$: At least four content-bearing words
        (length $> 4$, filtering stopwords by length heuristic) from the context
        appear in the response.
  \item \textbf{Bigram overlap} $> 1$: At least two consecutive-word pairs from
        the context appear in the response. Bigrams are a stronger signal because
        they capture phrasal matches rather than individual word coincidences.
\end{itemize}

A warning is emitted when the response is not grounded and the context contains
meaningful content (\texttt{contextWords.size > 10}).

The \texttt{groundingDetail} field in the assessment provides
\texttt{wordOverlap} and \texttt{bigramOverlap} counts for downstream
diagnostics.

% ============================================================================
\section{HanaDB Vector Store (Python)}
\label{sec:hana-db-python}
% ============================================================================

\noindent\textbf{Source:}
\texttt{src/data/langchain-integration-for-sap-hana-cloud-main/langchain\_hana/vectorstores/hana\_db.py}.

The \texttt{HanaDB} class extends LangChain's \texttt{VectorStore} base class to
provide SAP HANA Cloud Vector Engine integration.

\subsection{Distance Strategies}

Two distance functions are supported (lines 38--41):

\begin{lstlisting}[language=Python,caption={Distance function mapping --- lines 38--41}]
HANA_DISTANCE_FUNCTION: dict = {
    DistanceStrategy.COSINE:
        ("COSINE_SIMILARITY", "DESC"),
    DistanceStrategy.EUCLIDEAN_DISTANCE:
        ("L2DISTANCE", "ASC"),
}
\end{lstlisting}

\subsection{Internal vs.\ External Embeddings}

The class supports two embedding modes (lines 114--130):

\begin{lstlisting}[language=Python,caption={Embedding mode selection --- lines 114--130}]
def set_embedding(self, embedding: Embeddings) -> None:
    self.embedding = embedding
    if isinstance(embedding, HanaInternalEmbeddings):
        self.use_internal_embeddings = True
        self.internal_embedding_model_id = \
            embedding.get_model_id()
        self.internal_embedding_remote_source = \
            embedding.get_remote_source()
        self._validate_internal_embedding_function()
    else:
        self.use_internal_embeddings = False
        self.internal_embedding_model_id = ""
\end{lstlisting}

\begin{itemize}
  \item \textbf{Internal embeddings:} Use HANA's built-in \texttt{VECTOR\_EMBEDDING}
        SQL function. The embedding is computed server-side, eliminating network
        round-trips.
  \item \textbf{External embeddings:} Use a LangChain \texttt{Embeddings} instance
        (e.g., OpenAI, Cohere). Vectors are computed client-side and sent to HANA
        as binary payloads.
\end{itemize}

\subsection{Similarity Search with Score and Vector}

The core search method (lines 834--917) constructs a SQL query with optional
metadata filtering and keyword search projection:

\begin{lstlisting}[language=Python,caption={Core similarity search --- lines 879--896}]
sql_str = (
    f"{metadata_projection} "
    f"SELECT TOP {k}"
    f'  "{self.content_column}", '
    f'  "{self.metadata_column}", '
    f'  "{self.vector_column}", '
    f'  {distance_func_name}("{self.vector_column}", '
    f"  {embedding_expr}) AS CS "
    f"FROM {from_clause}"
)
parameters = []
if vector_embedding_params:
    parameters += vector_embedding_params
where_clause, where_parameters = \
    CreateWhereClause(self)(filter)
if where_clause:
    sql_str += f" {where_clause}"
    parameters += where_parameters
sql_str += (f" order by CS "
    f"{HANA_DISTANCE_FUNCTION[self.distance_strategy][1]}")
\end{lstlisting}

\subsection{Keyword Search Projection}

When a filter uses the \texttt{\$contains} operator on metadata fields, a CTE
(Common Table Expression) is generated to project JSON metadata values into
columns (lines 965--994):

\begin{lstlisting}[language=Python,caption={Metadata projection --- lines 983--994}]
metadata_columns = [
    (f"JSON_VALUE({self.metadata_column}, "
     f"'$.{HanaDB._sanitize_name(col)}') "
     f'AS "{HanaDB._sanitize_name(col)}"')
    for col in projected_metadata_columns
]
return (
    f"WITH {INTERMEDIATE_TABLE_NAME} AS ("
    f"SELECT *, {', '.join(metadata_columns)} "
    f"FROM \"{self.table_name}\")"
)
\end{lstlisting}

\subsection{Vector Serialization}

Embedding vectors are serialized in binary format for efficient HANA ingestion
(lines 364--375):

\begin{lstlisting}[language=Python,caption={Binary serialization --- lines 364--375}]
def _serialize_binary_format(self,
    values: list[float]) -> bytes:
    if self.vector_column_type == "HALF_VECTOR":
        return struct.pack(
            f"<I{len(values)}e", len(values), *values)
    elif self.vector_column_type == "REAL_VECTOR":
        return struct.pack(
            f"<I{len(values)}f", len(values), *values)
\end{lstlisting}

Two formats are supported:
\begin{itemize}
  \item \textbf{REAL\_VECTOR:} 4-byte IEEE 754 single-precision floats (standard
        FVECS format).
  \item \textbf{HALF\_VECTOR:} 2-byte IEEE 754 half-precision floats, reducing
        storage by 50\% with minimal accuracy loss for typical embedding models.
\end{itemize}

% ============================================================================
\section{Graph-RAG with KùzuDB}
\label{sec:graph-rag}
% ============================================================================

Two KùzuDB graph stores provide relationship context for the RAG pipeline.

\subsection{HANA LangChain Graph Store}

\noindent\textbf{Source:}
\texttt{src/data/langchain-integration-for-sap-hana-cloud-main/mcp\_server/kuzu\_store.py}.

\subsubsection{Schema Definition}

The schema (lines 54--86) defines three node types and three edge types:

\begin{lstlisting}[language=Python,caption={HANA KùzuDB schema --- lines 54--86}]
ddl = [
    # Node tables
    """CREATE NODE TABLE IF NOT EXISTS HanaVectorStore (
        storeId        STRING,
        tableName      STRING,
        embeddingModel STRING,
        schema         STRING,
        PRIMARY KEY (storeId)
    )""",
    """CREATE NODE TABLE IF NOT EXISTS EmbeddingDeployment (
        deploymentId  STRING,
        modelName     STRING,
        resourceGroup STRING,
        status        STRING,
        PRIMARY KEY (deploymentId)
    )""",
    """CREATE NODE TABLE IF NOT EXISTS HanaSchema (
        schemaId       STRING,
        schemaName     STRING,
        classification STRING,
        PRIMARY KEY (schemaId)
    )""",
    # Relationship tables
    """CREATE REL TABLE IF NOT EXISTS USES_DEPLOYMENT (
        FROM HanaVectorStore TO EmbeddingDeployment
    )""",
    """CREATE REL TABLE IF NOT EXISTS LIVES_IN (
        FROM HanaVectorStore TO HanaSchema
    )""",
    """CREATE REL TABLE IF NOT EXISTS RELATED_SCHEMA (
        FROM HanaSchema TO HanaSchema
    )""",
]
\end{lstlisting}

\begin{center}
\begin{tabular}{lll}
\toprule
\textbf{Edge Type} & \textbf{From} & \textbf{To} \\
\midrule
USES\_DEPLOYMENT & HanaVectorStore & EmbeddingDeployment \\
LIVES\_IN & HanaVectorStore & HanaSchema \\
RELATED\_SCHEMA & HanaSchema & HanaSchema \\
\bottomrule
\end{tabular}
\end{center}

\subsubsection{Graph Context Enrichment}

The \texttt{get\_store\_context()} method (lines 196--206) retrieves deployment
metadata linked to a vector store:

\begin{lstlisting}[language=Python,caption={get\_store\_context() --- lines 196--206}]
def get_store_context(self, store_id: str) -> list[dict]:
    """Return embedding deployments connected to a
    vector store."""
    return self.run_query(
        "MATCH (s:HanaVectorStore {storeId: $id})"
        "-[:USES_DEPLOYMENT]->"
        "(d:EmbeddingDeployment) "
        "RETURN d.deploymentId AS deploymentId, "
        "d.modelName AS modelName, "
        "d.resourceGroup AS resourceGroup, "
        "d.status AS status, "
        "'uses_deployment' AS relation",
        {"id": store_id},
    )
\end{lstlisting}

This graph context is attached to the RAG response in the MCP server's
\texttt{\_handle\_langchain\_rag\_chain()} method (server.py, lines 724--731):

\begin{lstlisting}[language=Python,caption={Graph context attachment --- server.py lines 724--731}]
if _kuzu_store_factory is not None:
    try:
        store = _kuzu_store_factory()
        if store.available():
            ctx = store.get_store_context(table_name)
            if ctx:
                base["graphContext"] = ctx
    except Exception:
        pass
\end{lstlisting}

\subsection{OData Vocabularies Graph Store}

\noindent\textbf{Source:}
\texttt{src/data/odata-vocabularies-main/mcp\_server/kuzu\_store.py}.

\subsubsection{Schema Definition}

The OData vocabulary store uses a distinct schema (lines 51--86):

\begin{lstlisting}[language=Python,caption={OData KùzuDB schema --- lines 51--86}]
ddl = [
    """CREATE NODE TABLE IF NOT EXISTS ODataVocabulary (
        vocabId    STRING,
        name       STRING,
        namespace  STRING,
        alias      STRING,
        status     STRING,
        PRIMARY KEY (vocabId)
    )""",
    """CREATE NODE TABLE IF NOT EXISTS VocabularyTerm (
        termId     STRING,
        name       STRING,
        type       STRING,
        appliesTo  STRING,
        description STRING,
        PRIMARY KEY (termId)
    )""",
    """CREATE NODE TABLE IF NOT EXISTS AnnotationTarget (
        targetId     STRING,
        entityType   STRING,
        namespace    STRING,
        usedInVocab  STRING,
        PRIMARY KEY (targetId)
    )""",
    """CREATE REL TABLE IF NOT EXISTS DEFINES_TERM (
        FROM ODataVocabulary TO VocabularyTerm
    )""",
    """CREATE REL TABLE IF NOT EXISTS ANNOTATES (
        FROM AnnotationTarget TO VocabularyTerm
    )""",
    """CREATE REL TABLE IF NOT EXISTS RELATED_VOCAB (
        FROM ODataVocabulary TO ODataVocabulary
    )""",
]
\end{lstlisting}

\begin{center}
\begin{tabular}{lll}
\toprule
\textbf{Edge Type} & \textbf{From} & \textbf{To} \\
\midrule
DEFINES\_TERM & ODataVocabulary & VocabularyTerm \\
ANNOTATES & AnnotationTarget & VocabularyTerm \\
RELATED\_VOCAB & ODataVocabulary & ODataVocabulary \\
\bottomrule
\end{tabular}
\end{center}

Both KùzuDB stores share the same operational pattern:
\begin{itemize}
  \item Graceful degradation when the \texttt{kuzu} package is absent.
  \item Singleton instantiation via \texttt{get\_kuzu\_store()}.
  \item Idempotent schema creation with \texttt{IF NOT EXISTS}.
  \item Parameterized Cypher queries preventing injection.
  \item \texttt{MERGE}-based upserts for idempotent data ingestion.
\end{itemize}

% ============================================================================
\section{Observability Spans}
\label{sec:observability}
% ============================================================================

\noindent\textbf{Source:}
\texttt{src/data/cap-llm-plugin-main/srv/cap-llm-plugin.ts},
\texttt{getRagResponseWithConfig()} at lines 449--596.

The RAG pipeline emits detailed OpenTelemetry span attributes for operational
monitoring. The following table summarizes the key span attributes:

\begin{longtable}{lll}
\toprule
\textbf{Span Attribute} & \textbf{Type} & \textbf{Description} \\
\midrule
\endhead
\texttt{rag.min\_score} & float & Minimum similarity score in filtered results \\
\texttt{rag.max\_score} & float & Maximum similarity score in filtered results \\
\texttt{rag.avg\_score} & float & Mean similarity score across filtered results \\
\texttt{rag.context\_chars} & int & Characters consumed from the context budget \\
\texttt{rag.context\_budget} & int & Total available context budget (characters) \\
\texttt{rag.documents\_used} & int & Number of documents included in the prompt \\
\texttt{rag.response.estimated\_tokens} & int & Estimated response token count ($\lceil \text{len}/4 \rceil$) \\
\texttt{rag.response.grounded} & bool & Whether the response is grounded in context \\
\bottomrule
\end{longtable}

These attributes are set at the following code locations:

\begin{lstlisting}[language=Java,caption={Score statistics --- lines 496--500}]
if (filteredResults.length > 0) {
  const scores = filteredResults.map(r => r.SCORE);
  span.setAttribute("rag.min_score",
    Math.min(...scores));
  span.setAttribute("rag.max_score",
    Math.max(...scores));
  span.setAttribute("rag.avg_score",
    scores.reduce((a, b) => a + b, 0) / scores.length);
}
\end{lstlisting}

\begin{lstlisting}[language=Java,caption={Context budget spans --- lines 532--534}]
span.setAttribute("rag.context_chars",
  maxContextChars - contextBudget);
span.setAttribute("rag.context_budget",
  maxContextChars);
span.setAttribute("rag.documents_used",
  selectedContent.length);
\end{lstlisting}

\begin{lstlisting}[language=Java,caption={Response quality spans --- lines 573--574}]
span.setAttribute("rag.response.estimated_tokens",
  quality.estimatedTokens);
span.setAttribute("rag.response.grounded",
  quality.usedContext);
\end{lstlisting}

Additional span events emitted during the pipeline:
\begin{itemize}
  \item \texttt{embedding\_generated} --- with \texttt{embedding.dimensions}.
  \item \texttt{similarity\_results\_filtered} --- with pre/post filter counts and
        threshold value.
  \item \texttt{rag\_no\_relevant\_context} --- when zero documents pass the
        similarity threshold.
  \item \texttt{chat\_completion\_received} --- after the LLM responds.
  \item \texttt{quality\_warnings} --- when the grounding check detects issues.
\end{itemize}

The similarity search method (\texttt{similaritySearch()}) also emits its own
spans at lines 665--669:
\begin{lstlisting}[language=Java,caption={Similarity search spans --- lines 665--669}]
const span = getTracer().startSpan(
  "cap-llm-plugin.similaritySearch");
span.setAttribute("db.hana.table", tableName);
span.setAttribute("db.hana.algo", algoName);
span.setAttribute("db.hana.top_k", topK);
span.setAttribute("db.hana.embedding_dims",
  embedding?.length ?? 0);
\end{lstlisting}

% ============================================================================
\section{System Architecture Summary}
\label{sec:architecture}
% ============================================================================

The embedding and RAG system follows a layered architecture:

\begin{enumerate}
  \item \textbf{Text Processing Layer:} Sentence-aware recursive splitting
        (Section~\ref{sec:recursive-split}) with paragraph $\to$ sentence $\to$
        word $\to$ character fallback hierarchy.

  \item \textbf{Embedding Layer:} Model-aware embedding generation
        (Section~\ref{sec:embedding-gen}) with dimension validation
        (Section~\ref{sec:model-registry}) and support for 15 models across
        4 providers.

  \item \textbf{Storage Layer:} Dual vector store implementations---Python
        HanaDB (Section~\ref{sec:hana-db-python}) with internal/external
        embeddings, and TypeScript CAPLLMPlugin
        (Section~\ref{sec:similarity-search}) with SDK/SQL dual path.

  \item \textbf{Graph Layer:} KùzuDB graph-RAG stores
        (Section~\ref{sec:graph-rag}) providing entity relationship context
        for vector stores, deployments, schemas, and OData vocabularies.

  \item \textbf{Quality Layer:} Score threshold filtering, context window
        budgeting, zero-result fallback, structured context injection,
        and response grounding assessment
        (Sections~\ref{sec:rag-quality}--\ref{sec:quality-assessment}).

  \item \textbf{Observability Layer:} Comprehensive OpenTelemetry spans
        (Section~\ref{sec:observability}) covering retrieval scores, context
        budgets, document usage, and response grounding.
\end{enumerate}

% ============================================================================
\section{File Reference Index}
\label{sec:file-index}
% ============================================================================

\begin{longtable}{p{0.55\textwidth}p{0.15\textwidth}p{0.25\textwidth}}
\toprule
\textbf{File Path} & \textbf{Lines} & \textbf{Section} \\
\midrule
\endhead
\texttt{src/data/langchain-integration-for-sap-hana-cloud-main/mcp\_server/server.py}
  & 72--128 & \ref{sec:recursive-split} Recursive Split \\
  & 412--424 & \ref{sec:mcp-splitter} MCP Tool Schema \\
  & 785--802 & \ref{sec:mcp-splitter} Split Handler \\
  & 724--731 & \ref{sec:graph-rag} Graph Context \\
\midrule
\texttt{src/data/cap-llm-plugin-main/lib/validation-utils.js}
  & 19--29 & \ref{sec:similarity-search} SQL Validation \\
  & 54--66 & \ref{sec:similarity-search} Vector Validation \\
  & 73--94 & \ref{sec:model-registry} Model Registry \\
  & 104--112 & \ref{sec:model-registry} Dim Validation \\
  & 122--183 & \ref{sec:quality-assessment} Quality Assessment \\
\midrule
\texttt{src/data/cap-llm-plugin-main/srv/cap-llm-plugin.ts}
  & 244--296 & \ref{sec:embedding-gen} Embedding Gen \\
  & 436--597 & \ref{sec:rag-quality} RAG Pipeline \\
  & 657--762 & \ref{sec:similarity-search} Similarity Search \\
\midrule
\texttt{src/data/langchain-integration-for-sap-hana-cloud-main/langchain\_hana/vectorstores/hana\_db.py}
  & 38--41 & \ref{sec:hana-db-python} Distance Strategies \\
  & 114--130 & \ref{sec:hana-db-python} Embedding Modes \\
  & 364--375 & \ref{sec:hana-db-python} Serialization \\
  & 834--917 & \ref{sec:hana-db-python} Search Query \\
  & 965--994 & \ref{sec:hana-db-python} Keyword Projection \\
\midrule
\texttt{src/data/langchain-integration-for-sap-hana-cloud-main/mcp\_server/kuzu\_store.py}
  & 54--86 & \ref{sec:graph-rag} HANA Schema \\
  & 196--206 & \ref{sec:graph-rag} Store Context \\
\midrule
\texttt{src/data/odata-vocabularies-main/mcp\_server/kuzu\_store.py}
  & 51--86 & \ref{sec:graph-rag} OData Schema \\
\bottomrule
\end{longtable}

\end{document}
